% Isolated after the CP33 canonical freeze. Complete proof; audits pending.
\subsection{The degree-one deck slice bounds the residual module}

Work geometrically over an algebraically closed field of characteristic
two. Let \(\pi:A\to B\) be an etale isogeny of odd degree \(n\),
where \(B\) is a nonordinary principally polarized abelian surface.
Put \(M=\mathcal O_B(3\Theta_B)\), \(L=\pi^*M\),
\(r=9n-1\), and let \(V\) be the geometric generic relative
Lang space for the complete system \(|L|\).
Let \(G\subseteq\operatorname{GL}(V)\) be its geometric image.
The canonical proper submodule \(K\) is the fixed space of the
normal closure of any full hyperplane radical, as in
Proposition~\ref{prop:root-full-radical-parabolic-form}.

\begin{lemma}[A block slice and a regular wall hyperplane]
\label{lem:DE-Lang-block-slice-bound}
\status{Candidate; complete direct proof supplied for independent audit.}
Let \(G\) be as above, and suppose a subgroup \(G_0\subseteq G\)
preserves a decomposition \(V=V_0\oplus V_1\). Suppose its
projection to \(\operatorname{GL}(V_0)\) is the full general
linear group, and \(G\) contains a full radical whose fixed
hyperplane is \(H=H_0\oplus V_1\), with \(H_0\) a hyperplane
of \(V_0\). Then \(K\subseteq V_1\).
\end{lemma}
\begin{proof}
The exact group-structure theorem makes \(K\) the unique maximal
proper \(G\)-submodule, independently of the initial full radical.
In particular it is fixed by the indicated full radical and lies
in \(H\). For every \(g\in G_0\), invariance of \(K\) gives
\(K\subseteq gH\). A defining functional of \(H\) is a nonzero
functional on \(V_0\), extended by zero on \(V_1\).
Its \(G_0\)-orbit spans \(V_0^*\), since the projection is the
full general linear group. The intersection of its kernels is
therefore precisely \(V_1\), proving the assertion.
\end{proof}

\begin{proposition}[A uniform eight-dimensional quotient]
\label{prop:DE-Lang-deck-slice-quotient-eight}
\status{Candidate; complete direct proof supplied for independent audit.}
For the odd cubic pullback above there is a nonempty open of the
degree-one slice \(\mathbf P(H^0(B,M))\subset\mathbf P(H^0(A,L))\)
on which the relative Lang local system splits over \(\mathbf F_2\)
as
\[
                 V=V_0\oplus V_{\mathrm{nt}},\qquad
                 \dim V_0=8,
 \tag{DE-DS1}
\]
where \(V_0\) is the pulled-back degree-one relative local system
and \(V_{\mathrm{nt}}\) is the nontrivial deck-character part.
After choosing a geometric base point in this open and the associated
identification with the generic space, one has
\[
 K\subseteq V_{\mathrm{nt}},\qquad
 \dim(V/K)\geq8,\qquad \dim K\leq r-8.
 \tag{DE-DS2}
\]
Consequently the unique main Lang component has degree at least
\(2^r-2^{r-8}\). For every fixed nonordinary principal surface
and every odd \(d\), the same conclusions hold with
\(r=9d^2-1\) and an outgoing maximal-isotropic quotient.
\end{proposition}
\begin{proof}
\paragraph{The actual splitting on a dense slice open.}
Let \(D=\ker\pi\), an odd constant finite group. Its averaging
idempotent is defined over \(\mathbf F_2\): since \(|D|\) is odd,
\[
                   e_D=\sum_{g\in D}g,
                   \qquad e_D^2=e_D.
\]
The deck action commutes with coherent Frobenius and hence acts on
the actual \(\mathbf F_2\)-Lang sheaf. The decomposition into the
image and kernel of \(e_D\) therefore exists over \(\mathbf F_2\),
without choosing individual character eigenlines over that field.

On coherent sections over the algebraically closed field, the
decomposition is the one in (DE-CI5). The trivial block is the
degree-one Euler bundle of \(B\); every nontrivial block has
dimension nine and its Frobenius character cycles are bijective
where their multiplication determinants are nonzero. The relative
trivial block is exactly that of \(B\), since odd etale pullback
identifies the constant ambient cohomology as well. Thus
Proposition~\ref{prop:DE-common-Cartier-incidence} supplies a nonempty
open \(S\) of the degree-one smooth, sharp section space on which
all these blocks are bijective and (DE-DS1) holds. In particular
this is a splitting of actual local systems, not only a decomposition
of a root-removed matrix.

Let \(S_1\) be the degree-one smooth, sharp open before removing
the additional nontrivial-block determinants. The map
\(\pi_1(S)\to\pi_1(S_1)\) is surjective: a connected finite etale
cover of the normal integral scheme \(S_1\) is normal and integral,
and remains connected on any nonempty open. The proved degree-one
fixed-surface theorem gives image \(\operatorname{GL}_8(\mathbf F_2)\)
on \(V_0\) over \(S_1\), including when \(B\) is a product.
Consequently the slice image \(G_0\subseteq G\) preserves (DE-DS1)
and projects onto that full group. No surjectivity from the slice
fundamental group to the full ambient fundamental group is asserted.

\paragraph{The ambient radical has a slice-supported normal.}
Use the fixed theta coordinates of
Lemma~\ref{lem:odd-cover-common-Fermat-coordinates} on this slice.
Choose a geometric point \(x\) on its integral Fermat cubic where
the central Pfaffian cubic is nonzero, the central first and stable
ranks are seven, and all nontrivial multiplication and cross blocks
are invertible. Such points form a nonempty open by
Proposition~\ref{prop:odd-cover-common-rank-wall-nonempty} and
Lemma~\ref{lem:nonordinary-uniform-Q-stable-loss-candidate}.

Let \(X=U_{\max}\) be the ambient common maximal-rank open and
\(H_{\max}\subset X\) its wall. This wall is smooth and reduced;
its first and stable relative ranks are \(r-1\). Its intersection
with the slice at \(x\) is the reduced Fermat wall. In particular
there is a transverse punctured slice neighborhood inside the
Frobenius-bijective locus. On the slice, the regular wall Lang
space is precisely
\[
                    H_0\oplus V_{\mathrm{nt}},
                    \qquad \dim H_0=7:
 \tag{DE-DS3}
\]
the trivial block has stable rank seven, while every other character
cycle remains bijective.

We spell out the comparison with the \emph{ambient} radical. Pass
to a strictly henselian local neighborhood \(R\) of \(x\) in \(X\).
The Lang equations on the regular Euler bundle have Jacobian
\(-I\) in the fibre variables. Every one of the \(2^{r-1}\)
regular Lang vectors at \(x\) therefore has a unique etale lift
over \(R\). The lifts form a constant rank-\((r-1)\)
\(\mathbf F_2\)-subspace \(\mathcal H\) on the punctured
neighborhood: addition and all relations follow from uniqueness
of lifting. At the generic point of the wall through \(x\),
these are all the regular Lang vectors, since its stable rank is
still \(r-1\).

The ambient inertia there is the full hyperplane radical of
Theorem~\ref{thm:odd-cover-full-hyperplane-inertia}. Its fixed
hyperplane is consequently \(\mathcal H\). Its inclusion in the
global geometric image makes the canonical \(K\) lie in
\(\mathcal H\) at a geometric generic point of this punctured
neighborhood. The inclusion persists on that connected neighborhood,
as an inclusion of finite etale local systems. Restricting to the
punctured slice compares exactly the unique lifts of the same
regular points. It identifies \(\mathcal H\) with (DE-DS3).
This argument specializes only the regular rank-\((r-1)\) part;
it makes no specialization isomorphism for the full rank-\(r\)
generic fibre across a rank drop.

Thus, at a base point of \(S\) reached from this local slice
neighborhood, the ambient full radical has fixed hyperplane
\(H_0\oplus V_{\mathrm{nt}}\). Its defining covector is nonzero
on \(V_0\) and zero on \(V_{\mathrm{nt}}\). The preceding lemma,
with the actual slice image \(G_0\), gives (DE-DS2).
Changing the base point only transports the statement with the
specified slice splitting.

Finally Lemma~\ref{lem:DE-Lang-main-orbit} identifies the main
component degree with \(2^r-2^{\dim K}\), giving the stated bound.
For \(|3d\Theta|\), use the exact outgoing etale quotient of
degree \(d^2\) in
Corollary~\ref{cor:fixed-ppav-all-odd-maximal-dual-Cartier}.
It supplies precisely the present setup and proves the last sentence.
The case \(n=1\), or \(d=1\), has \(V_{\mathrm{nt}}=0\) and
agrees with the already proved degree-one full image.
\end{proof}

\begin{remark}[Scope of the slice bound]
\label{rem:DE-Lang-deck-slice-scope}
\status{Candidate scope; no comparison of different deck slices.}
This is a bound on the already canonical exceptional module.
It neither proves that \(K=0\) in higher degree nor identifies
the relative position of distinct deck slices after a common
specialization. The full group on the central summand is used
only as a projection of the slice image; the image on the other
summand need not be independent. No new computation or external
classification enters the proof. The constant regular wall part,+rather than an entire rank-dropping fibre, is the comparison object.
\end{remark}
